\subsection{Pointed equivariantizations and an affine group}
Four further models are strict unitary equivariantizations of untwisted
pointed categories. Their group automorphisms are
\[
\begin{array}{lll}
R8\_2\_244:& A=C_4^2,&(x,y)\mapsto(-y,x-y),\quad G=C_3;\\
R8\_2\_289:& A=C_2^4,&(x,y,z,w)\mapsto(y,x+y,w,z+w),\quad G=C_3;\\
R8\_2\_143:& A=C_{17},&x\mapsto4x,\quad G=C_4;\\
R8\_2\_179:& A=C_2^4,&(x,y,z,w)\mapsto(w,x+w,y+w,z+w),\quad G=C_5.
\end{array}
\]
In each case the displayed generator has the stated order and acts
freely off zero. The zero orbit supplies the ordinary characters of
$G$, and each nonzero orbit has trivial stabilizer. Tensor products
are calculated by orbit addition and projection to stabilizer
characters, giving the exact full source tensors. Strict permutation
functors with identity tensorators supply coherent unitary actions.

The ring $R8\_2\_237$ is
$K_0(\Rep(\operatorname{A\Gamma L}(1,8)))$. This group of order
$168$ has elements $(v,a,t)$ acting on $\mathbb F_8$ as
$x\mapsto v+a x^{2^t}$. Its simple dimensions are $(1^3,3^2,7^3)$.
The three dimension-seven characters have the exact expression
\[
\chi_s(v,a,t)=\zeta_3^{st}
\sum_{\substack{r\in\mathbb F_8^\times\\ar^{2^t-1}=1}}
 (-1)^{\operatorname{Tr}_{\mathbb F_8/\mathbb F_2}(v/r)}.
\]
Together with the characters inflated from
$\mathbb F_8^\times\rtimes C_3$, they give an exact character table
in $\Q(\zeta_{21})$. Orthogonality and all tensor multiplicities are
checked independently by finite group sums.

\subsection{Bosonic quotients and the \texorpdfstring{$D_{14}$}{D14} model}
The ring $R8\_2\_557$ is the central bosonic quotient of
$PSU(2)_6\boxtimes PSU(2)_6$ by its diagonal invertible. Each factor
invertible is a fermion, so their product has self-braiding $+1$ and
trivial pointed associator. It generates a genuine central
$\Rep(C_2)$. All sixteen simples have free diagonal orbits; the eight
free-module orbits give the stated ring.

For $R8\_2\_313$, use instead the two unitary pseudo-$PSU(2)_6$
categories of \cite{EIP}, with the exact sourced associators discussed
above. Their central invertibles have self-braidings $+i$ and $-i$,
and each has associator $F_{ggg}=-1$. The diagonal product therefore
has both trivial self-braiding and trivial associator; its central
$\Rep(C_2)$ gives the same eight-orbit quotient construction.
For each rank-four factor the universal grading is trivial, so a
central lift of an invertible is unique; its central square is the unit.
This verifies the order-two condition as well as the associator and
self-braiding conditions required for the categorical quotient.

The ring $R8\_2\_855$ is the even $D_{14}$ category, or the local
module category of the $SU(2)_{24}$ algebra $1\oplus V_{24}$.
It is unitary modular \cite{MPSskein,MPSknots}. Its basis is
$f_0,f_2,f_4,f_6,f_8,f_{10},P,Q$, with $P,Q$ self-dual. Fold the bulk
$SU(2)_{24}$ products by $f_{2t}=f_{2(12-t)}$ and $f_{12}=P+Q$.
The remaining rules are
\[
\begin{gathered}
P f_{2j}=\sum_{k=6-j}^{5}f_{2k}
 +\begin{cases}P&j\text{ even},\\Q&j\text{ odd},\end{cases}
\qquad Q f_{2j}\text{ interchanges }P,Q,\\
P^2=P+f_0+f_4+f_8,\quad Q^2=Q+f_0+f_4+f_8,
\quad PQ=f_2+f_6+f_{10}.
\end{gathered}
\]
All these constructions preserve the unitary structure.

\subsection{A Morita construction over \texorpdfstring{$C_7$}{C7}}
The ring $R8\_2\_239$ is a Morita dual of an $S_3$ extension of
$\Vec_{C_7}$. In
$\operatorname{BrPic}(\Vec_{C_7})=O(C_7\oplus\widehat{C_7})$ take
$r=\operatorname{diag}(2,4)$ and let $s$ exchange the two coordinates.
They generate $S_3$. Its degree-three obstruction vanishes by coprime
cohomology with seven-primary coefficients; the next obstruction lies
in $H^4(S_3,\C^\times)=0$. Thus the actual extension exists
\cite[Theorem~1.3]{ENOExt}. Its $C_3$ restriction is pointed, on
$F_{21}=C_7\rtimes C_3$. The associator class there is inflated from
$C_3$: the order-three action kills the $H^3(C_7,\C^\times)$ invariant
part. Surjectivity of
$H^3(S_3,\C^\times)\to H^3(C_3,\C^\times)$ permits an outer twist
that makes this restriction untwisted.

Take the separable algebra $B=\bigoplus_{h\in C_3}h$ in this
restriction. The category of $B$-bimodules is an actual Morita dual,
whose neutral category is $\Rep(F_{21})$. Write its simples as
$\alpha_i$ ($i\in\Z/3$), $X,Y$, and odd $m_j$ ($j\in\Z/3$).
The nontrivial rules, with all $m_j$ self-dual, are
\[
\begin{gathered}
X^2=X+2Y,\quad Y^2=2X+Y,\quad
XY=\sum_i\alpha_i+X+Y,\\
\alpha_i m_j=m_{i+j},\quad m_j\alpha_i=m_{j-i},\qquad
X m_j=Y m_j=m_jX=m_jY=\sum_km_k,\\
m_jm_k=\alpha_{j-k}+X+Y.
\end{gathered}
\]
They follow from the $S_3$ bimodule action and reciprocity, and match
all $512$ source coefficients. The extension and its Morita dual are
weakly group-theoretical, hence unitarizable
\cite[Theorem~5.20]{GHR}.

\subsection{Charge-conjugation defects of \texorpdfstring{$PSU(3)_4$}{PSU(3)4}}
The ring $R8\_2\_290$ is a unitary $C_2$ crossed extension of the
rank-five modular category $\mathcal B=PSU(3)_4$. Charge conjugation
is an actual braided symmetry
\cite[Corollary~2.7]{EdieSU3}. The invertible group of $\mathcal B$
is trivial, so the degree-three gauging obstruction vanishes, and
$H^4(C_2,U(1))=0$ removes the remaining obstruction. The unitary
gauging theorem \cite[Sections~3.2--3.4]{Gauging} therefore supplies
an actual unitary extension. Its odd rank is three, the number of
charge-conjugation-fixed neutral simples \cite[Proposition~1.1]{Bischoff}.

For completeness, its tensor is forced by a small exact module
calculation. In the neutral basis
$(0,0),(3,0),(1,1),(0,3),(2,2)$ the nonselfdual generator $X$ has
regular characteristic polynomial
$(t^2+t+2)(t^3-2t^2-t+1)$, computed by quantum Weyl alternants in
$\Q(\zeta_{21})$. Its odd action $A$ has the cubic factor as
characteristic polynomial, since its Perron root has degree three.
The commuting neutral action and the dual-transpose relation make
$A$ normal. Its three roots are real, so $A$ is symmetric.
Nonnegative integer entries, $\operatorname{Tr}(A)=2$ and
$\operatorname{Tr}(A^2)=6$ give, up to relabelling,
\[
A=\begin{pmatrix}0&0&1\\0&1&1\\1&1&1\end{pmatrix}.
\]
The five neutral actions are
$I,A,A^2-I,A,A^2-A$. Distinct odd dimensions force each odd object
to be self-dual; reciprocity then determines every odd product.
The independent constructor matches all $512$ census coefficients.

\subsection{A self-dual Hopf extension in dimension thirty-two}
For $R8\_2\_575$ the neutral Hopf algebra is
$H=H_{C:\sigma_1}$ of dimension sixteen from \cite{Kashina}. A
symmetric nondegenerate Hopf pairing supplies the extra data of the
Davydov--Runkel construction \cite[Theorem~2.5]{DavydovRunkel}.
Here is enough explicit information to fix the construction and its
normalization. Write $h(p,q,j)$ for its basis, with $p$ modulo four
and $q,j$ modulo two. Put $\zeta=\exp(2\pi i/16)$ and $i=\zeta^4$.
The pairing is
\[
P(h(p,q,j),h(r,s,k))
 =\delta_{s,j}\delta_{k,q}\,i^{-3rp}\zeta^{jq}/4.
\]
The supplementary Hopf presentation and checker verify nondegeneracy,
symmetry and every pairing identity exactly in $\Q(\zeta)$.
Davydov--Runkel use $H^\vee=(H^*)^{\rm op,cop}$; consequently their
self-duality map is $\Gamma=S\circ P^{-1}$, rather than the
unmodified inverse pairing. The right cointegral
$\lambda(h(p,q,j))=\delta_{j,0}/2$ satisfies
$\Gamma(\lambda)=4\Lambda$ and
$\lambda(\Gamma(\lambda))=1$, where $\Lambda$ is the normalized
integral. The distinguished grouplike is one, $S^2=1$, and $\Gamma$
is symmetric, so all conditions of the cited theorem hold.
Its rigidity criterion \cite[Proposition~3.25]{DavydovRunkel} gives
an actual fusion category $\Rep(H)\oplus\Vec$. Its neutral ring is
$\Rep(D_{16})$; for the odd object $M$ of dimension four,
\[
X M=M X=\FPdim(X)M,\qquad
M^2=\sum_{X\in\Irr(H)}\FPdim(X)X.
\]
The dimension-$32$ solvability theorem and
\cite[Theorem~5.20]{GHR} make this fusion category unitarizable.
The complete model tensor, Hopf identities, and source comparison are
separate exact checks.

\subsection{The dual Haagerup--Izumi \texorpdfstring{$C_7$}{C7} category}
\label{sec:hi7dual}
The ring $R8\_2\_868$ is the dual even part of the unitary
Haagerup--Izumi $C_7$ subfactor. Its existence and center are supplied
by \cite[Theorem~5(d)]{EvansGannon}. Put
$d=(7+\sqrt{53})/2$, so $d^2=1+7d$.
Izumi's dual $Q$-system analysis
\cite[Proposition~7.4 and Theorem~7.5]{Izumi3n} gives rank eight,
with dimensions $1,d,(d-1)^3,(d+1)^3$.

We verify the canonical-algebra hypothesis needed to use the full
fusion rules in \cite[Theorem~7]{EvansGannon}; matching modular data
alone would not suffice. In its published double, the twist-one
simples are $1$, $b$ of dimension $d^2$, and six $A_i$ of dimension
$1+d^2$. The nonzero $C_{53}$ sector has no twist one, and the other
twist-one labels are precisely those with zero second coordinate.
For the actual dual category $\mathcal D$, its canonical Lagrangian
algebra $L=I_{\mathcal D}(1)$ is connected, has trivial twist and
dimension $7(1+d^2)$. Therefore
\[
L=1+m_b b+\sum_{i=1}^6m_iA_i.
\]
Writing $T=\sum_i m_i$, rational and irrational coefficients in the
dimension identity give $m_b+T=7$ and $m_b+2T=13$, hence
$m_b=1,T=6$. Adjunction gives
$\dim\operatorname{End}(L)=\rank(\mathcal D)=8$, so
$\sum_i m_i^2=6$. Thus every $m_i=1$.
The actual canonical algebra has exactly the underlying object
required by the cited theorem; no assumption about a guessed algebra
multiplication is made.

For a direct exact constructor, use basis
$1,a,b_1,b_2,b_3,c_1,c_2,c_3$. Put $B=\sum_i b_i$, $C=\sum_i c_i$,
$b_0=c_0=0$, and $b_{-i}=b_i,c_{-i}=c_i$ for indices modulo seven.
The rules are
\[
\begin{aligned}
a^2&=1+a+B+C,& ab_i&=a+B-b_i+C,\\
ac_i&=a+B+c_i+C,& b_ic_j&=a+B+C,\\
b_ib_j&=\delta_{ij}1+(1-\delta_{ij})a+B-b_{i+j}-b_{i-j}+C,\\
c_ic_j&=\delta_{ij}1+(1+\delta_{ij})a+B+c_{i+j}+c_{i-j}+C.
\end{aligned}
\]
They match every source coefficient; all $4096$ associativity
equations and all dimension products in $\Q(d)$ are checked.
Both even parts of the cited finite-depth $C^*$ subfactor are unitary.

\begin{center}\small
\begin{longtable}{@{}lp{0.65\textwidth}@{}}\toprule
Model & Model-to-census permutation\\\midrule
$R8\_2\_143$ & $(0,1,7,2,3,5,6,4)$\\
$R8\_2\_179$ & $(0,1,4,3,2,5,6,7)$\\
$R8\_2\_237$ & $(0,1,2,5,6,7,4,3)$\\
$R8\_2\_239$ & $(0,1,2,3,4,5,7,6)$\\
$R8\_2\_244$ & $(0,1,2,3,7,4,5,6)$\\
$R8\_2\_289$ & $(0,1,2,3,4,5,6,7)$\\
$R8\_2\_290$ & $(0,1,6,2,7,3,4,5)$\\
$R8\_2\_313$ & $(0,7,1,2,3,4,6,5)$\\
$R8\_2\_557$ & $(0,1,2,7,3,6,5,4)$\\
$R8\_2\_575$ & $(0,5,1,2,3,7,4,6)$\\
$R8\_2\_855$ & $(0,1,7,4,5,6,2,3)$\\
$R8\_2\_868$ & $(0,3,1,2,4,5,7,6)$\\\bottomrule
\end{longtable}
\end{center}
The model bases are fixed by the independent constructors in the
supplement. Each permutation is checked on the entire ordered tensor;
it is not inferred from a character table or dimension vector.
